\scanpage{50}
% archive foliation: 49
groupe des autom. de $(X,\Delta)$ qui fixe $\infty$ et $-1$, donc \inserted{des autom.} de la forme $z\longmapsto az+b$ avec $-a+b=-1$ i.e.
\[
z\longmapsto\deleted{\unsure{a(z-1)+1}=az+(1-a)}\quad az+(a-1)
\]
est plus petit, en fait si l'autom. envisagé n'est pas l'identité, il doit induire une permutation sans pts fixes de $(0,1,i,-i)$, qui est donc d'ordre $2$ ou quatre, donc on a $a^4=1$ \deleted{ou $a^4=1$}.
\deleted{i.e. $a\in\unsure{\{1,i,-i\}|\{0,1,i,-i\}}$}
mais faisant $z=0$ on doit avoir
\[
a-1\in\{1,i,-i\}\qquad\text{i.e.}\quad a\in(2,1+i,1-i)
\]
faisant $z=1$ on trouve
\[
2a-1\in\{0,i,-i\}\qquad\text{i.e.}\quad a\in\left\{\frac12,\frac12(i+1),\frac12(1-i)\right\}
\]
\margin{donc $a^4=\pm4$ qui ne peut être $=1$ qu'en car. $3$ et $5$. \unsure{faux}, $a^4=-1$}
\margin{donc $a^4=1/4$ qui n'est pas ég. à $1$ (car. on est en car. $\ne2$) [\unsure{faux}.]}
faisant $z=i$, on trouve
\[
a\underbrace{(i+1)}_{1:\left(\frac{1-i}{2}\right)}-1\in\{0,1,-i\}
\qquad\text{i.e.}\quad
 a\in\left\{\deleted{\unsure{\frac1{i+1}}}\inserted{\unsure{\frac{1-i}{2}}},\deleted{\unsure{\frac2{i+1}}}\inserted{1-i},-i\right\}
\]
et de $=$, faisant $z=-i$, on trouve
\[
a\in\left\{\frac{1+i}{2},1+i,i\right\}.
\]
Donc l'\deleted{\unsure{a}}intersection des stabilisateurs de $\infty,-1$ est $\{1\}$, et par raison d'homogénéité il en est de $=$ pour toute partie à deux éléments de $\Delta$, distincte des $\Delta_1,\Delta_2,\Delta_3$. Donc on voit que la structure octaédrale $\Delta=\Delta_1\amalg\Delta_2\amalg\Delta_3$ est \unsure{reconnue|préservée} par $\Gamma$, et uniquement déterminée par le diviseur $\Delta$. Donc $\Gamma\subset\underline{\operatorname{Aut}}_{\mathrm{oct}}(\Delta)$, Pour montrer qu'on a isomorphisme, il suffit de noter que le groupe des autom. \inserted{directs} de la structure octaédrale est engendré par les trois sous-groupes cycliques d'ordre $4$ \deleted{\unsure{tous}} qu'il contient, correspondant à $\Delta_1,\Delta_2,\Delta_3$ -- or ceux-ci sont contenus dans $\Gamma$.
\margin{\unsure{et donc|et même} $\Gamma\subset\operatorname{Aut}_{\mathrm{oct}}^+(\Delta)$ \unsure{$=$}}
